Distributed acoustic sensing (DAS) converts an optical fibre into thousands of vibration sensors, and its value for perimeter security depends on the recognition of the events that those sensors record. The deep neural networks that lead public DAS benchmarks generally require graphics processing units (GPUs), are difficult to interpret, and are usually evaluated on random or chronological splits of overlapping windows. This thesis aimed to analyse the computational footprint of DAS event classification and to establish what makes a pipeline lightweight, to develop a multi-domain machine learning pipeline for the task, and to evaluate that pipeline against published deep learning baselines. Following the Design Science Research Methodology, the pipeline was designed on the nine-class Tomasov perimeter security dataset and applied without modification to the six-class Cao dataset, which was recorded with an unrelated sensing system. Each window is represented by $1{,}002$ named descriptors drawn from the time, frequency, time--frequency, and spatial domains, a set that a consensus of eight feature selectors can reduce, and seven classifiers were compared under six strategies for handling class imbalance. The novelty of the work is methodological rather than algorithmic. Train--test disjointness is asserted in code before any model is trained, and every model is profiled for size, latency, and training cost on the same embedded board. An audit of the released protocols showed that, under a random split, every Tomasov test window shared raw samples with a training window, and that 423 of the 441 Cao recording sessions appeared in both published folders. Blocked partitions with guard bands and partitions grouped by session eliminated both forms of leakage. On these partitions, the fusion of the four domains outperformed every single domain, and the selected LightGBM classifier attained macro-$F_1$ scores of $81.98$ on Tomasov and $98.90$ on Cao. These scores exceed those of the convolutional neural networks published with the two datasets by $16.43$ and $8.01$ points, and those of five graph and transformer architectures published since 2024. Consensus selection removed approximately $95\%$ of the features at a cost of $8.39$ and $5.81$ points, whereas the choice of balancing strategy had little effect. Changing the partition alone altered the accuracy of an otherwise identical network by up to $26.27$ points, and removing the leakage reduced the $F_1$ score of the \emph{openclose} class from a published $0.72$ to $0.044$. Generalisation beyond the training recordings was weaker. Accuracy averaged $63.85\%$ on unseen recordings of four Tomasov classes and fell from $99.06\%$ to $86.81\%$ on unseen Cao acquisition dates, and, on the four event types common to the two datasets, a model trained on one sensing system performed below chance on the other. On an NVIDIA Jetson AGX Orin, the selected model occupies $0.83$--$3.42$\,MB, is trained on the CPU in $41$--$1{,}357$\,s, and classifies a window in under $0.2$\,ms, whereas the neural baselines occupy up to $132.6$\,MB and require up to ten and a half hours of GPU training. Feature extraction, at $122$--$145$\,ms per window, remains the most computationally expensive stage. The results demonstrate that a lightweight, interpretable pipeline is competitive with recent deep architectures, and that such comparisons are meaningful only when the evaluation is grouped by recording and verified in code.
